\section{Introduction}
\label{sec:intro}

Derived categories of line operators of $3$-dimensional topological quantum field theories (TQFTs) form braided monoidal, i.e.~$\mathbb{E}_2$-monoidal $\infty$-categories, and often arise via
deformation quantization of $2$-shifted symplectic stacks.
As such, they are candidates for determining $4$-dimensional theories.
More explicitly, if an extended $3$-dimensional theory assigns some object $B$ to the point,
then its line operators are its value on the circle, which is the $\EE_2$-algebra
$\mathcal{Z}_1\left(B\right)$.
The object $B$ itself can be viewed relative to $\mathcal{Z}_1\left(B\right)$, hence we can ask when it determines a boundary theory for the
associated $4$-dimensional theory.

In this article we consider a zoo of different examples of $B$ which are suitable sheaf-theoretic categories and as such inherently derived.
We investigate when $\mathcal{Z}_1\left(B\right)$ is {\em invertible}, so that
the $4$-dimensional theory is invertible and is an anomaly of the $3$-dimensional theory. Concretely, $\mathcal{Z}_1\left(B\right)$ is an object in the higher Morita category $\umor_2\left(\Pre\right)$, and we discuss invertibility as such. By \cite{BJSS, SSS} (Theorem \ref{thm:inv_E2}), this holds if and only if it is $3$-dualizable,
non-degenerate, factorizable and cofactorizable; these conditions are phrased using factorization homology.

The examples are motivated by Rozansky--Witten theory, which takes as input a holomorphic
symplectic space, or more generally a $2$-shifted symplectic stack.
Such a space is locally modeled on a symplectic vector space, so the global theory should
be determined by the linear theory and the action of its symmetries (providing
the transition data).
We expect this action to be anomalous, with the anomaly obstructing the gluing.
This is why we focus on the seemingly simple linear case, and why invertibility is the
question of interest: As explained above, the line operators define an anomaly precisely
when they are invertible.
Quasi-coherent sheaves serve as a warm-up for shifted $\cD$-modules, which provide an avatar
for the line operators of linear Rozansky--Witten theory.
The same framework includes Chern--Simons theory, whose input is the $2$-shifted symplectic
stack $BG$.
When $G$ is a torus, this is a global example which is understood, also in the derived
setting, through its quantum group.

\subsection*{Quasi-coherent sheaves}
As a warm-up, we revisit quasi-coherent sheaves and the formalism of Ben-Zvi, Francis and Nadler in \cite{BZFN,BZFN:app}.
Using their formulas for factorization homology we rephrase the invertibility conditions geometrically.
It follows immediately that a commutative algebra is invertible only if its unit is an
equivalence, so that $\left(\QC\left(X\right),\otimes\right)$ is invertible only if $X$ is a point.
Instead, a basic source of invertible
examples is convolution categories on loop stacks $\left(\QC\left(\cL X\right),*\right)$, which are the Drinfeld centers of $\left(\QC\left(X\right),\otimes\right)$.
In order to calculate their factorization homologies, we pass through 
a Morita equivalence between 
sheaves on the loop stack $\cL X= X\times_{X\times X} X$ with convolution and sheaves on the base $X\times X$ with the symmetric monoidal tensor product. 
For loop stacks, under the hypothesis of diagonal descent, we upgrade the Morita equivalence of \cite{BZFN:app} to $\umor_2(\Pre)$, using the criterion of \cite[Theorem 2.26]{BJSS}.


\begin{introthm}[Proposition \ref{prop:inv_symm}, Theorems \ref{thm:inv_rel} and \ref{thm:trivial_QC}]
\label{thm:intro_A}
\leavevmode
\begin{enumerate}
\item For $d\geq1$ and $\coef$ as in Section \ref{sec:mor}, a commutative algebra in $\coef$ is
invertible in $\umor_d\left(\coef\right)$ if and only if its unit is an equivalence. In particular, if $X$ is a quasi-compact stack with affine
diagonal and a $k$-point, then $\left(\QC\left(X\right),\otimes\right)$ is invertible in
$\umor_d\left(\Pre\right)$ only if $X\simeq\pt$.
\item If $X$ satisfies diagonal descent, e.g.\ $X=BG$ or the formal completion $\widehat V_0$ of a
vector space at the origin, then $\left(\QC\left(\cL X\right),*\right)$ is Morita trivial,
trivialized by $\QC\left(X\right)$.
\end{enumerate}
\end{introthm}

Part (2) is the derived-geometric analogue of the Witt triviality of Drinfeld centers of fusion
categories \cite{DMNO}. Diagonal descent is restrictive, however: It fails for any
positive-dimensional scheme (Remark \ref{rmk:diag_descent_fails}), and in general we do not
expect $\QC\left(\cL X\right)$ to be invertible.

\subsection*{Shifted D-modules}
Next, we turn to our original motivation for this article, namely studying the line operators of Rozansky--Witten theory.
Choosing the target $T^*\left[2\right]X$, these are closely related to the
deformation quantization $\DMod^2\left(X\right)$ of the $2$-shifted symplectic stack
$T^*\left[2\right]X$ \cite{Ber:Sph, HoLi} (Section \ref{sec:dmod}).
We show that it is typically not invertible.
To test the invertibility conditions, we compute the higher centers of $\DMod^2\left(X\right)$ using the
excision theorem of \cite{HoLi}, together with rigidity of $\DMod^n\left(X\right)$ for $X$ bounded (Proposition \ref{prop:rigid}), which is observed for $n=1$ in
\cite{Ber:center} and used in \cite[Section 5.2]{Ber:Sph}. Rigidity also gives $3$-dualizability of $\DMod^2\left(X\right)$ for $X$
smooth (Corollary \ref{cor:3dual}).
The symmetric monoidal category $\DMod\left(X\right)$ of ordinary $\cD$-modules is central in
$\DMod^2\left(X\right)$, so $\DMod^2\left(X\right)$ sheafifies over the de Rham stack of $X$
(Remark \ref{rmk:DMod_center}).
More precisely, for $X=BG$, $\DMod^2\left(X\right)$ is non-degenerate exactly when $G$ is
unipotent (Theorem \ref{thm:centers} and Proposition \ref{prop:BG}). Even then, as for an odd vector
space $X=V\left[1\right]$, invertibility can still fail through (co)factorizability.
We therefore expect $\DMod^2\left(X\right)$ to be invertible only when $X$ is a point.

\begin{introthm}[Theorems \ref{thm:centers} and \ref{thm:BW}, Proposition \ref{prop:BG} and Corollary \ref{cor:BG_not_inv}]
\label{thm:intro_B}
\leavevmode
\begin{enumerate}
\item For a smooth stack $X$ as in Section \ref{sec:dmod} and $j=0,1,2$ (with $X\to\cL X$ bounded if $j=2$, e.g.\ $X=BG$), the $\EE_j$-center of $\DMod^2\left(X\right)$ is its
factorization homology over the $j$-sphere:
\begin{equation*}
\cZ_j\left(\DMod^2\left(X\right)\right)\simeq\int_{S^j\times\RR^{2-j}}\DMod^2\left(X\right)
\simeq\DMod^{2-j}\left(X^{S^j}\right) \ .
\end{equation*}
\item\label{thm:intro_B-example} For a non-zero vector space $V$, $\DMod^2\left(V\left[1\right]\right)$ is $3$-dualizable and
non-degenerate, but neither factorizable nor cofactorizable. In particular, it is not invertible in
$\umor_2\left(\Pre\right)$.
\item If $G$ is a non-trivial affine algebraic group which is either not unipotent or abelian, then
$\DMod^2\left(BG\right)$ is not invertible.
\end{enumerate}
\end{introthm}

For finite braided tensor categories over an algebraically closed field of characteristic zero,
non-degeneracy, factorizability and cofactorizability are equivalent to each other and to invertibility \cite{BJSS}. By Theorem \ref{thm:intro_B}(\ref{thm:intro_B-example}), in the
derived setting non-degeneracy implies neither of the other two, even for $3$-dualizable
$\EE_2$-algebras. The reason is visible already in Theorem \ref{thm:intro_B}(1): For $X=V\left[1\right]$,
the mapping stacks $X^{S^0}$ and $X^{S^2}$ are purely odd, whereas
$X^{S^1}\simeq V\times V\left[1\right]$ has an even direction.

\subsection*{Equivariantization}
Next we turn to symmetries. A symmetry of a $3$-dimensional theory by an algebraic group $G$
induces an action of $G$ on its category of line operators (Section \ref{sec:gauging_lines}).
Given such an action on a category of line operators $A$ by
$\mathbb{E}_2$-monoidal functors, one can form its \emph{equivariantization} $A^G$, which
is an $\mathbb{E}_2$-algebra among $\Rep\left(G\right)$-modules, i.e.\ an object of
$\umor_2\left(\Mod_{\Rep\left(G\right)}\left(\Pre\right)\right)$.
This is the (de)equivariantization of \cite{DGNO} for finite groups and of \cite{AG:small}
for quantum groups, formulated using $1$-affineness of $BG$ \cite{G:ShvCat}.

\begin{introthm}[Proposition \ref{prop:equiv}, Remark \ref{rmk:inv_descent} and Lemma \ref{lem:geometric}]
\label{thm:intro_equiv}
Let $G$ be an affine algebraic group of finite type over a field of characteristic zero, acting on
an $\EE_2$-algebra $A$ in $\Pre$ by $\EE_2$-functors.
\begin{enumerate}
\item Equivariantization gives an equivalence
\begin{equation*}
\umor_2\left(\ShvCat\left(BG\right)\right)\simeq\umor_2\left(\Mod_{\Rep\left(G\right)}\left(\Pre\right)\right)
\ ,
\end{equation*}
and $A^G$ is invertible in the latter if and only if $A$ is invertible in $\umor_2\left(\Pre\right)$.
\item If the action comes from an action of $G$ on a perfect stack $X$ satisfying diagonal
descent, with $X/G$ perfect and
$A=\left(\QC\left(\cL X\right),*\right)$, then $\QC\left(X/G\right)$ trivializes
$A^G\simeq\QC\left(\cL X/G\right)$ over $\Rep\left(G\right)$.
\end{enumerate}
\end{introthm}

When $A$ is invertible, we define the \emph{anomaly} of the action as the class
$\left[A^G\right]-\left[\Rep\left(G\right)\otimes A\right]$ in the Witt group of
$\Rep\left(G\right)$-linear $\EE_2$-algebras (Definition \ref{defn:anomaly}). 
By (2), symmetries of the $2$-shifted cotangent bundle $T^*\left[2\right]X$ which are induced
by symmetries of $X$ itself are anomaly-free, at least when $X$ satisfies diagonal descent and
the line operators are $\QC\left(\cL X\right)$. This suggests that anomalies come from
symmetries of $T^*\left[2\right]X$ which do not come from $X$.
Two such symmetries motivate this paper. The first is the group of linear symmetries of
$T^*\left[2\right]\left(V\left[1\right]\right)\simeq\left(V\dsum V^\vee\right)\left[1\right]$, the
target of linear Rozansky--Witten theory, preserving its shifted symplectic structure; this is
larger than the group $\mathrm{GL}\left(V\right)$ of symmetries of $V\left[1\right]$.
The second is the quantum Frobenius symmetry of the small quantum group, to which we turn
next.

\subsection*{Quantum groups of tori}
\label{sec:intro_CS}

The line operators of Chern--Simons theory for a complex reductive algebraic group at a
root of unity are expected to form the $\EE_2$-category $\Rep\left(G_q\right)$ of
representations of the big quantum group.
Quantum Frobenius gives a central $\EE_2$-functor $\Rep\left(H\right)\to\Rep\left(G_q\right)$
for a dual group $H$, and $\Rep\left(G_q\right)$ is the equivariantization of an action of
$H$ on a version $A$ of the small quantum group \cite{AG:small, Neg:frob, Neg:log}; these
results are in the abelian setting (Remark \ref{rmk:abelian}).
For a torus, $A$ is a non-degenerate pointed braided fusion category, and for fusion categories
nothing is lost in passing to derived categories. So for tori we can work in the derived setting.

\begin{introthm}[Theorem \ref{thm:T_invertible}]
\label{thm:intro_C}
Let $k=\CC$, let $T$ be a torus with character lattice $\Lambda$, let $Q\colon\Lambda\to k^\times$ be a quadratic
form whose double braiding has radical $\Lambda_0$ of finite index with $Q|_{\Lambda_0}=1$, let
$\Rep\left(T_q\right)$ be the braided category of $\Lambda$-graded complexes determined by $Q$, and let
$H\ceqq\Hom\left(\Lambda_0,\GG_m\right)$. Then $\Rep\left(T_q\right)$ is the equivariantization of
an action of $H$ on the non-degenerate pointed braided fusion category $A=\Vect_{\Lambda/\Lambda_0}$,
and it is invertible in
\begin{equation*}
\umor_2\left(\Mod_{\Rep\left(H\right)}\left(\Pre\right)\right) \, .
\end{equation*}
\end{introthm}

Here $\Rep\left(T_q\right)$ and $A$ are the categories of representations of the big and small
quantum groups of $T$. The weight modules of Lusztig's divided-power quantum group of $T$ are
$\Lambda$-graded, with braiding given by $Q$, and $\Rep\left(H\right)=\Vect_{\Lambda_0}$ is the
M\"uger center of $\Rep\left(T_q\right)$. When $\Lambda_0$ is the character lattice of the target
of quantum Frobenius, its inclusion is quantum Frobenius, and $A=\Vect_{\Lambda/\Lambda_0}$ is the
category of representations of the small quantum group of $T$, which is the group algebra of the
finite group $\Hom\left(\Lambda/\Lambda_0,\GG_m\right)$ \cite{AG:small}. For instance, for
$T=\GG_m$ and $Q\left(a\right)=q^{a^2}$, with $q$ a primitive $n$-th root of unity and $n$ odd,
$\Lambda_0=n\ZZ$, quantum Frobenius is the isogeny $z\mapsto z^n$, and the small quantum group is
$k\left[K\right]/\left(K^n-1\right)$.

So the anomaly of the action of $H$ on $A$ is defined (Definition \ref{defn:anomaly}); we
do not determine here whether it vanishes. Detecting it, by transgression to character
stacks of surfaces, will be treated in forthcoming work.

\subsection*{Outlook: an analytic flavor}
In Section \ref{sec:aggregated} we briefly discuss an analytic setting, aggregated vector spaces, in
which more Witt classes should become trivial. We show that, given suitable classes attached to
real quadratic spaces, the image of a non-degenerate pointed braided fusion category there only depends on its
central charge modulo $8$ (Proposition \ref{prop:Witt_reduction}).
This analytic context, and various TQFTs defined in this setting, will be the subject of
forthcoming work with Pelle Steffens.

\subsection*{Organization} 

Section \ref{sec:Morita} recalls higher Morita categories and the criteria for
dualizability and invertibility. Section \ref{sec:sheaves} recalls quasi-coherent sheaves
on stacks and their factorization homology, and Section \ref{sec:inv_sheaves} studies
their invertibility. 
We introduce and prove our results concerning shifted $\cD$-modules in Section
\ref{sec:dmod}. Section \ref{sec:gauging} treats equivariantization and anomalies,
including quantum groups of tori and our expectations in an analytic setting.

\subsection*{Acknowledgments}

We thank David Ben-Zvi, Dan Freed, Pavel Safronov, Germ\'an Stefanich and Constantin Teleman
for many inspiring conversations. Special thanks go to Pelle Steffens for many useful
discussions, in particular on the analytic setting of Section \ref{sec:aggregated}.

Both authors were supported by the Simons Collaboration on Global Categorical Symmetries
(1013836 and 8528-03).
Both authors were supported by the Deutsche Forschungsgemeinschaft (DFG, German Research
Foundation) through the Collaborative Research Centers SFB 1085 Higher invariants -
224262486 and SFB 1624 Higher structures, moduli spaces and integrability - 506632645.

